\documentclass[12pt,a4paper]{article}

%% ── Packages ────────────────────────────────────────────────────────────────
\usepackage[utf8]{inputenc}
\usepackage[T1]{fontenc}
\usepackage{lmodern}
\usepackage{amsmath, amssymb}
\usepackage{graphicx}
\usepackage{booktabs}
\usepackage{siunitx}
\usepackage{hyperref}
\usepackage{geometry}
\usepackage{caption}
\usepackage{subcaption}
\usepackage{cleveref}

\geometry{margin=2.5cm}
\sisetup{per-mode=symbol}

%% ── Title ───────────────────────────────────────────────────────────────────
\title{Design and Simulation of Terahertz Coplanar Waveguide Resonators}
\author{Abdullah Alkabbawi}
\date{\today}

\begin{document}

\maketitle
\tableofcontents
\newpage

%% ── 1. Introduction ─────────────────────────────────────────────────────────
\section{Introduction}
\label{sec:intro}

This report presents the design and simulation of terahertz (THz)
coplanar waveguide (CPW) resonators using Ansys Lumerical FDTD\@.
The aim is to understand how resonator geometry controls the resonance
frequency, transmission response, bandwidth, and loaded quality factor.

Two coupling methods are investigated:
\begin{enumerate}
  \item Simple capacitive gap coupling.
  \item Interdigital capacitor (IDC) coupling.
\end{enumerate}

%% ── 2. CPW Resonator Design ─────────────────────────────────────────────────
\section{CPW Resonator Design}
\label{sec:design}

\subsection{Half-wavelength resonator}
\label{subsec:half-wave}

The initial cavity length is obtained from the half-wavelength relation:
\begin{equation}
  L_{\mathrm{cav}} = \frac{c}{2\,n_{\mathrm{eff}}\,f_0},
  \label{eq:halfwave}
\end{equation}
where $f_0 = \SI{0.8}{\tera\hertz}$, $n_{\mathrm{eff}} = 2.65$,
and $c = \SI{3e8}{\metre\per\second}$.
This gives $L_{\mathrm{cav}} \approx \SI{70.75}{\micro\metre}$.

\subsection{CPW dimensions}
\label{subsec:cpw-dims}

\begin{table}[h]
  \centering
  \caption{Baseline CPW dimensions.}
  \label{tab:cpw-dims}
  \begin{tabular}{llr}
    \toprule
    Parameter              & Symbol             & Value \\
    \midrule
    Centre conductor width & $s$                & \SI{10}{\micro\metre} \\
    CPW slot gap           & $w$                & \SIrange{6.6}{7.5}{\micro\metre} \\
    Target frequency       & $f_0$              & \SI{0.8}{\tera\hertz} \\
    Initial $n_\text{eff}$ & $n_{\mathrm{eff}}$ & 2.65 \\
    Initial cavity length  & $L_{\mathrm{cav}}$ & \SI{70.75}{\micro\metre} \\
    \bottomrule
  \end{tabular}
\end{table}

%% ── 3. Gap Coupling ─────────────────────────────────────────────────────────
\section{Simple Gap-Coupled Resonator}
\label{sec:gap}

\subsection{Method}
\label{subsec:gap-method}

Simple capacitive gaps of width $g$ separate the feedlines from the
resonant cavity.  The loaded quality factor is:
\begin{equation}
  Q_L = \frac{f_0}{\Delta f},
  \label{eq:QL}
\end{equation}
where $\Delta f$ is the full width at half maximum (FWHM) of the
Lorentzian transmission peak.

\subsection{Results}
\label{subsec:gap-results}

\begin{table}[h]
  \centering
  \caption{Fitted parameters from the gap-coupled resonator simulations.}
  \label{tab:gap-results}
  \begin{tabular}{rrrr}
    \toprule
    Coupling gap & $f_0$ (\si{\tera\hertz}) & $\Delta f$ (\si{\giga\hertz}) & $Q_L$ \\
    \midrule
    \SI{5.0}{\micro\metre}   & 0.6675 & 49.25  & 13.55 \\
    \SI{3.0}{\micro\metre}   & 0.6670 & 58.78  & 11.35 \\
    \SI{1.0}{\micro\metre}   & 0.6454 & 92.09  & 7.01  \\
    \SI{0.8}{\micro\metre}   & 0.6399 & 101.03 & 6.33  \\
    \bottomrule
  \end{tabular}
\end{table}

%% ── 4. End Correction ───────────────────────────────────────────────────────
\section{Effective Index and End Correction}
\label{sec:endcorr}

A corrected resonance model is:
\begin{equation}
  f_0 = \frac{c}{2\,n_{\mathrm{eff}}\,(L_{\mathrm{cav}} + \Delta L)},
  \label{eq:corrected}
\end{equation}
where $\Delta L$ accounts for fringing fields and capacitive loading.
Fitting to the length-sweep data gives $n_{\mathrm{eff}} \approx 2.68$
and $\Delta L \approx \SI{12.38}{\micro\metre}$.

%% ── 5. IDC Coupling ─────────────────────────────────────────────────────────
\section{Interdigital Capacitor Coupling}
\label{sec:idc}

\subsection{IDC geometry}
\label{subsec:idc-geom}

\begin{table}[h]
  \centering
  \caption{IDC geometry parameters.}
  \label{tab:idc-params}
  \begin{tabular}{llr}
    \toprule
    Parameter      & Symbol                 & Value \\
    \midrule
    IDC length     & $L_{\mathrm{IDC}}$     & \SI{12}{\micro\metre} \\
    Finger width   & $w_f$                  & \SI{0.6}{\micro\metre} \\
    Finger gap     & $g_f$                  & \SI{0.45}{\micro\metre} \\
    Tip gap        & $g_{\mathrm{tip}}$     & \SI{0.5}{\micro\metre} \\
    Fingers/side   & $N$                    & 4 \\
    Total fingers  & $2N$                   & 8 \\
    \bottomrule
  \end{tabular}
\end{table}

%% ── 6. Data Processing ──────────────────────────────────────────────────────
\section{Data Processing}
\label{sec:data}

The exported quantity is $|S_{21}|^2$, converted to decibels by:
\begin{equation}
  S_{21,\mathrm{dB}} = 10\log_{10}\!\left(|S_{21}|^2\right).
  \label{eq:dB}
\end{equation}

Reference-normalised transmission is:
\begin{equation}
  S_{21,\mathrm{norm,dB}} =
    10\log_{10}\!\left(\frac{|S_{21}|^2}{|S_{21,\mathrm{ref}}|^2}\right).
  \label{eq:norm}
\end{equation}

Quality factors are related by:
\begin{equation}
  \frac{1}{Q_L} = \frac{1}{Q_i} + \frac{1}{Q_c}.
  \label{eq:Q-relation}
\end{equation}

%% ── 7. Conclusions ──────────────────────────────────────────────────────────
\section{Conclusions}
\label{sec:conclusions}

% TODO: fill in after IDC length sweep is complete.

\appendix

%% ── Bibliography ────────────────────────────────────────────────────────────
\bibliographystyle{ieeetr}
\bibliography{references}

\end{document}
